\section{Introduction}

Building energy-flexibility research considers how operational demand can be adjusted while respecting service constraints~\cite{Chen2018,Junker2018}. Occupancy is an important contextual signal for that task: it can identify intervals in which a building zone may be empty and, therefore, potentially suitable for a carefully constrained operational action. Occupancy modeling has long been studied for building energy management~\cite{Erickson2014}, and recent reviews show a broad forecasting literature spanning data acquisition, estimation, prediction, and evaluation~\cite{Jin2021,CaballeroPena2024,Li2024}. Operational recommendations add three requirements to forecast discrimination: timely issuance, useful duration, and an acceptable observed conflict rate.

In plain terms, the analysis has two stages. First, three models estimate how strongly each future 15-min interval resembles an Empty interval. Their scores are combined in a \emph{hybrid}: a weighted average of model outputs. Second, a decision rule converts sufficiently high scores into recommendations that continue for at least one hour. Conflict rate measures the share of recommended intervals later labelled occupied by the cameras. A load-opportunity proxy records processed HVAC-plus-lighting load during recommended intervals later labelled Empty.

Occupancy-prediction-based cooling-control studies and HVAC-control evaluations assess modeled or field operating interventions~\cite{Peng2018,EsrafilianNajafabadi2022,Lin2023}. This study measures an \emph{offline safe load opportunity}: the processed HVAC-plus-lighting proxy in intervals recommended by a fixed rule and subsequently labelled Empty by the cameras. The statistic screens candidate intervals for operational study. Energy-savings estimation requires a control action, a counterfactual, and explicit comfort and equipment constraints.

The repository's left-labelled records determine forecast timing. A row labelled 00:00 summarizes $[00{:}00,00{:}15)$ and becomes available at 00:15, so the study evaluates post-bin next-record forecasts. Section~VI describes the source-timestamp and imputation provenance available in the cleaned release.

The research question is: \emph{for one office-building time series, how do validation-selected Empty-class scores trade off forecast discrimination, subsequently observed camera-label conflict, and offline safe load opportunity when recommendations must form at least one-hour stable windows?} Forecast discrimination measures how well the models rank future Empty intervals across all overlapping rolling 24-h forecasts. Policy evaluation converts those scores into recommendations on non-overlapping midnight-labelled daily horizons, counting each target interval once.

The paper makes three contributions. First, it documents a post-bin chronological protocol that maps 24-h Empty-class scores to stable-window recommendations and reserves controllable-load data for downstream evaluation. Second, it evaluates a validation-selected hybrid of a historical schedule baseline, LightGBM, and a compact Transformer. The weights optimize validation forecast discrimination; the operating threshold maximizes the validation load-opportunity proxy under a 10\% empirical conflict-rate limit. Third, it reports the uncertainty behind the point estimates. The hybrid has slightly higher AUPRC than the two baselines and zero observed camera-label conflicts in the held-out policy period. Paired test contrasts cross zero, and validation perturbations reveal several competitive weights and thresholds.

Section~II positions the evaluation against occupancy forecasting and control research. Sections~III and~IV define the data, timing, policy, and experimental protocol. Section~V reports the canonical results, and Sections~VI--VII discuss their limits and implications. Appendix~A records the status of the expanded decision-aware and window-aware searches.
